\documentclass[../../main.tex]{subfiles}
\begin{document}

{\bf [解]}
題意より
\begin{align}
 n\in\mathbb{Z}_{\ge3} \label{eq:1}\\
 x,y,z\in\mathbb{N} \label{eq:2}
\end{align}
である．このもとで与えられた不等式
\begin{align}
\begin{cases}\displaystyle
x+y+z=n \\
x\le y+z \\
y\le z+x \\
z\le x+y
\end{cases}
\quad\cdots *_1 \label{eq:3}
\end{align}
を考える．

まず $z=k\in\mathbb{N}$と固定して\cref{eq:3}に代入して
\begin{align}
\begin{cases}\displaystyle
  x+y=n-k \\
  k\le x+y \\
  y\ge x-k \\
  y\le x+k
\end{cases}
\quad\cdots *_2 \label{eq:4}
\end{align}
となる．第一式より$y=n-k-x$だから，これを残りの3つの不等式に代入して
\begin{align}
 \begin{cases}
  k \le n-k \\
  x-k\le n-k-x\le x+k   
 \end{cases} \\
\therefore
\begin{cases}
   k \le \dfrac{n}{2} \\
  \dfrac{n}{2}-k\le x\le\dfrac{n}{2}\quad\cdots\text{④}
\end{cases} \label{eq:5} 
\end{align}
である．また，$y>0$より
\begin{align}
 x < n-k \label{eq:6}
\end{align}
である．

\subparagraph{$n=2p\ (p\in\mathbb{N})$の時（$p=2,3,\cdots$）}

この時\cref{eq:5}より$k=1,2,\cdots, p$をとる．
よって$z=k$と固定した時の$(x,y)$の数は，\cref{eq:5,eq:6}より
\begin{align}
 \begin{cases}\displaystyle
  k=1,2,\cdots,p-1\text{の時} & k+1\text{個} \\
  k=p\text{の時} & p-1\text{個}
 \end{cases}
\end{align}
である．よってこれらを足し合わせて求める数$S_n$は
\begin{align}
 S_{2p} 
   &=\sum_{k=1}^{p-1}(k+1)+(p-1) \\
   &=\left(\dfrac{p(p+1)}{2}-1\right)+(p-1) \\
   &=\dfrac{1}{2}p^2+\dfrac{3}{2}p-2 \\
   &=\dfrac{1}{8}n^2+\dfrac{3}{4}n-2\ \text{個}
\end{align}
である．

\subparagraph{$n=2p+1\ (p\in\mathbb{N})$の時（$p=1,2,3,\cdots$）}

この時は\cref{eq:5}より$k=1,2,\cdots,p$をとる．
よって$z=k$と固定した時の$(x,y)$の数は，\cref{eq:5,eq:6}より$k$個である．
よってこれらを足し合わせて求める数$S_n$は
\begin{align}
  S_{2p+1} 
   &=\sum_{k=1}^{p}k \\
   &=\dfrac{1}{2}p(p+1) \\
   &=\dfrac{1}{2}\cdot\dfrac{n-1}{2}\cdot\dfrac{n+1}{2}\ \text{個}
\end{align}
である．

以上二つの場合分けより求める数は
\begin{align}
\begin{cases}\displaystyle
\dfrac{1}{8}n^2+\dfrac{3}{4}n-2\text{個} & (n:\text{even}) \\[2mm]
\dfrac{1}{2}\cdot\dfrac{n-1}{2}\cdot\dfrac{n+1}{2}\text{個} & (n:\text{odd})
\end{cases}
\end{align}
である．

\newpage

\end{document}
